% =====================================================================
%  Speech de soutenance — version vérifiée
% =====================================================================
\documentclass[a4paper,11pt]{article}
% =====================================================================
%  Préambule commun aux documents de soutenance
% =====================================================================
\usepackage[T1]{fontenc}
\usepackage[utf8]{inputenc}
\usepackage[sfdefault]{carlito}
\usepackage[french]{babel}
\usepackage{csquotes}
\usepackage[a4paper,margin=2.2cm,top=2.5cm,bottom=2.2cm]{geometry}
\usepackage{microtype}
\usepackage[htt]{hyphenat}
\tolerance=2000
\emergencystretch=3em
\widowpenalty=10000
\clubpenalty=10000
\setlength{\parskip}{5pt plus 2pt minus 1pt}

\usepackage{xcolor}
\usepackage{colortbl}
\usepackage{array}
\usepackage{tabularx}
\usepackage{longtable}
\usepackage{enumitem}
\usepackage{ragged2e}
\usepackage{titlesec}
\usepackage{fancyhdr}
\usepackage{tcolorbox}
\usepackage[hidelinks]{hyperref}

\definecolor{btkblue}{HTML}{14507A}
\definecolor{btklight}{HTML}{1B6CA8}
\definecolor{btkpale}{HTML}{EAF3FA}
\definecolor{btkambre}{HTML}{9A6B12}
\definecolor{btkambrepale}{HTML}{FDF6E6}
\definecolor{btkvert}{HTML}{2E7D5B}
\definecolor{btkvertpale}{HTML}{EDF7F2}

\renewcommand{\tabularxcolumn}[1]{>{\RaggedRight\arraybackslash}p{#1}}
\newcolumntype{L}[1]{>{\RaggedRight\arraybackslash}p{#1}}
\renewcommand{\arraystretch}{1.25}

\titleformat{\section}{\Large\bfseries\color{btkblue}}{\thesection}{0.7em}{}
\titleformat{\subsection}{\large\bfseries\color{btklight}}{\thesubsection}{0.6em}{}
\titlespacing*{\section}{0pt}{2.4ex plus 1ex}{1.2ex}
\titlespacing*{\subsection}{0pt}{1.8ex plus .8ex}{0.8ex}

\pagestyle{fancy}
\fancyhf{}
\renewcommand{\headrulewidth}{0.4pt}
\renewcommand{\headrule}{\hbox to\headwidth{\color{btklight}\leaders\hrule height \headrulewidth\hfill}}
\fancyhead[R]{\small\color{btkblue}Oussema Korbosli --- PFE BTK}
\fancyfoot[C]{\small\thepage}

% Encadré « à savoir par cœur »
\newtcolorbox{coeur}[1][]{
  colback=btkpale, colframe=btklight, boxrule=0.9pt, arc=2pt,
  left=7pt, right=7pt, top=5pt, bottom=5pt,
  fonttitle=\bfseries\small, title={À savoir par c\oe{}ur}, #1}

% Encadré « note » : conseil de présentation
\newtcolorbox{note}[1][]{
  colback=btkvertpale, colframe=btkvert, boxrule=0.8pt, arc=2pt,
  left=7pt, right=7pt, top=4pt, bottom=4pt, #1}

% Encadré « chiffre corrigé »
\newtcolorbox{corrige}[1][]{
  colback=btkambrepale, colframe=btkambre, boxrule=0.9pt, arc=2pt,
  left=7pt, right=7pt, top=5pt, bottom=5pt,
  fonttitle=\bfseries\small, title={Chiffre corrigé}, #1}

% Question / réponse
\newcommand{\question}[1]{\subsection{#1}}
\newcommand{\cle}[1]{\par\vspace{2pt}\noindent
  \colorbox{btkpale}{\parbox{\dimexpr\linewidth-2\fboxsep}{%
  \textbf{Phrase à dire :} \emph{#1}}}\par\vspace{3pt}}

\fancyhead[L]{\small\color{btkblue}Speech de soutenance}

\begin{document}

\begin{center}
{\color{btkblue}\rule{\textwidth}{2pt}}\\[0.3cm]
{\huge\bfseries\color{btkblue} Speech de soutenance\par}
\vspace{0.2cm}
{\large Système intelligent de gestion des rôles et de pointage\\
avec tableau de bord décisionnel --- BTK\par}
\vspace{0.2cm}
{\color{btkblue}\rule{\textwidth}{2pt}}
\end{center}

\vspace{0.4cm}
\begin{corrige}[title={Version vérifiée --- ce qui a changé}]
Les chiffres de ce speech ont été recoupés avec le code, le datamart et le
rapport. \textbf{Cinq corrections} ont été faites, toutes sur la partie
segmentation et sur le modèle en étoile. Elles sont récapitulées en dernière
page. Les chiffres du réseau (45 agences, 997 employés, 29\,942 clients,
105\,000 comptes, 1,40 milliard de crédits, 253 millions d'épargne,
87,9\,\% de présence) ont été vérifiés et sont \textbf{exacts}.

\medskip
\textbf{Une diapositive a été ajoutée} : le \emph{diagramme de classes},
inséré en \textbf{slide 19}, juste après le diagramme de cas d'utilisation. La
présentation passe de 36 à \textbf{37 diapositives}, et \textbf{tout ce qui
suit est décalé d'un rang} --- le modèle en étoile devient la 20, la conclusion
la 35. Les numéros de ce document sont déjà à jour.
\end{corrige}

% =====================================================================
\section*{Slide 1 --- Page de garde}
\addcontentsline{toc}{section}{Slide 1}

Bonjour à tous. Je m'appelle Oussema Korbosli, et je vais vous présenter mon
projet de fin d'études intitulé \enquote{Système intelligent de gestion des
rôles et de pointage avec tableau de bord décisionnel}, réalisé au sein de la
BTK --- Banque Tuniso-Koweïtienne, sous l'encadrement de M. Melek Aziz Elmoula
comme maître de stage, et de M. Heni Abidi comme encadrant académique.

\section*{Slide 2 --- Plan}

Je vais structurer ma présentation en six parties : le contexte général, les
spécifications des besoins, la conception et l'architecture, le volet
décisionnel et sa réalisation, l'intelligence artificielle avec une
démonstration, et enfin les perspectives d'évolution.

\begin{note}
\small Le plan est déjà écrit sur la diapositive : ne pas le lire, faire une
simple liaison.
\end{note}

% =====================================================================
\section*{Partie 01 --- Contexte général (slides 3 à 10)}

\subsection*{Slide 4 --- Contexte général}

La BTK gère un réseau de 45 agences, avec des données éclatées entre plusieurs
fichiers et applications, alors que la direction doit piloter la performance et
l'assiduité sur des indicateurs fiables et à jour. Comme je le résume :
\enquote{Une donnée dispersée ne décide rien : c'est en la consolidant qu'elle
devient un instrument de pilotage.} L'objectif : centraliser la gestion du
réseau et transformer les données en indicateurs d'aide à la décision.

\subsection*{Slide 5 --- Présentation de l'organisme d'accueil}

J'ai effectué mon stage à la BTK, une banque universelle tunisienne créée le
25 février 1981 dans le cadre d'une convention de coopération avec le Koweït,
au département Data, du 16 février au 16 juin 2026, sous l'encadrement de
M. Melek Aziz Elmoula.

\subsection*{Slide 6 --- Le réseau BTK en chiffres}

Le réseau en chiffres : 45 agences, 997 employés, près de 30\,000 clients,
plus de 105\,000 comptes, 1,40 milliard de dinars de crédits, 253 millions
d'épargne, et 87,9\,\% de présence.

Ces chiffres sont issus des \textbf{données réelles de la base BTK}, restituées
par le rapport Power BI --- pas de données fictives.

\begin{corrige}
La phrase d'origine attribuait ces chiffres \enquote{au datamart produit par ma
chaîne ETL}. C'est imprécis et attaquable : le datamart de l'ETL agrège un
périmètre plus restreint (20\,979 comptes, 378 millions de crédits). Les
105\,000 comptes et le 1,40 milliard viennent de \texttt{B\_OBJECTIF} agrégée
sur \textbf{toutes les périodes} par Power BI. Les deux sont réels, mais ce ne
sont pas les mêmes agrégations : attribuez-les à Power BI, pas à l'ETL.
\end{corrige}

\subsection*{Slide 7 --- Problématiques}

Quatre problématiques en découlent : des données dispersées et consolidées
manuellement, un suivi de présence sans traçabilité fiable, aucune vision
consolidée du réseau, et des indicateurs indisponibles au moment de décider.

\subsection*{Slide 8 --- Étude de l'existant}

J'ai étudié trois solutions : Excel, risqué et sans traçabilité des opérations
sensibles ; un progiciel du marché, coûteux et peu adapté aux besoins
spécifiques de la BTK. J'ai donc retenu un développement sur mesure, adapté au
métier réel de la banque.

\subsection*{Slide 9 --- Solution proposée}

Ma solution a trois volets : une application web de gestion centralisée des
agences, employés, clients et objectifs, en Jakarta EE avec une base Oracle ;
une chaîne décisionnelle ETL vers un datamart en étoile restitué via Power BI ;
et un volet intelligence artificielle de segmentation des agences avec
scikit-learn. Une seule solution intégrée, de la saisie jusqu'à la décision.

\subsection*{Slide 10 --- Méthodologie adoptée : CRISP-DM}

J'ai adopté la méthodologie CRISP-DM, en six phases : compréhension métier,
compréhension des données, préparation des données, modélisation, évaluation,
et déploiement.

Pourquoi ce choix ? Mon projet est centré sur les données : Scrum structure bien
la production logicielle, mais ne couvre pas l'analyse et la valorisation des
données, alors que CRISP-DM est la méthodologie de référence des projets
décisionnels, et ses six phases épousent exactement le déroulement de mon
projet --- tout en gardant des pratiques agiles comme les itérations et les
validations régulières.

\begin{coeur}
La justification de CRISP-DM contre Scrum est une question quasi certaine.
Retenir l'argument : \emph{Scrum organise la production logicielle, CRISP-DM
organise le travail sur la donnée}.
\end{coeur}

% =====================================================================
\section*{Partie 02 --- Spécifications des besoins (slides 11 à 14)}

\subsection*{Slide 12 --- Identification des acteurs}

Trois acteurs : l'Administrateur gère les agences, employés et clients, valide
les demandes et supervise le décisionnel ; le Directeur commercial soumet les
demandes de création et valide les demandes de modification des employés ;
l'Utilisateur consulte ses données et effectue son pointage. Ces rôles sont
gérés via la table \texttt{SECURITY\_USERS}.

\subsection*{Slide 13 --- Besoins fonctionnels}

Six besoins fonctionnels : l'authentification et la gestion des rôles, la
gestion du référentiel --- agences, employés, clients ---, le suivi des
objectifs, le pointage avec ses statuts, les circuits de validation, et le
volet décisionnel.

\subsection*{Slide 14 --- Besoins non fonctionnels}

Six exigences non fonctionnelles : la sécurité, la fiabilité, la performance,
l'ergonomie, la maintenabilité, et la portabilité --- le standard Jakarta EE
rendant l'application compatible avec différents serveurs.

% =====================================================================
\section*{Partie 03 --- Conception et architecture (slides 15 à 20)}

\subsection*{Slide 16 --- Architecture logique en couches}

Une architecture logique en quatre couches, inspirée du MVC : Présentation,
Contrôle, Métier et Persistance, chacune indépendante et faiblement couplée aux
autres --- ce qui facilite la maintenance.

\begin{coeur}
Les quatre couches et leur rôle : c'est le premier point que le jury vérifie.
Présentation = les JSP ; Contrôle = les servlets ; Métier = les EJB ;
Persistance = JPA/Hibernate.
\end{coeur}

\subsection*{Slide 17 --- Architecture physique (3-tiers)}

Sur le plan physique, trois niveaux : le poste client, le serveur WildFly qui
héberge l'application, et le serveur Oracle --- avec Power BI et la chaîne ETL
Python connectés directement aux vues décisionnelles de la base.

\subsection*{Slide 18 --- Diagramme de cas d'utilisation global}

Ce diagramme illustre les interactions des trois acteurs.
\enquote{S'authentifier} est commun à tous via une relation d'inclusion.
L'administrateur concentre la gestion, le directeur commercial soumet et
consulte, l'utilisateur pointe et consulte ses données.

\subsection*{Slide 19 --- Diagramme de classes du modèle métier}

Ce diagramme présente le modèle métier : \textbf{treize classes}, regroupées en
\textbf{cinq paquetages fonctionnels} pour en faciliter la lecture --- le
référentiel du réseau, la gestion commerciale, la présence, les workflows de
demandes, et la sécurité et la traçabilité.

La classe \textbf{Utilisateur} occupe une position centrale : c'est elle que
référencent le pointage, les demandes, le journal d'audit et les notifications.
Les rattachements à l'agence et au gestionnaire sont portés par des \textbf{clés
étrangères} --- \texttt{SK\_AGENCE} et \texttt{SK\_GESTIONNAIRE} ; les entités de
demande, le pointage, le journal et les notifications référencent l'utilisateur
par une association \texttt{@ManyToOne}. Seule \texttt{DemandeAgence} entretient
une association \texttt{@ManyToMany} avec \texttt{Utilisateur}, puisqu'une
demande de création d'agence porte la liste des employés à y rattacher.

Le modèle ne comporte ni héritage ni composition : ce sont des associations,
implémentées soit par des clés étrangères, soit par des relations JPA
explicites. Enfin, la classe \texttt{Pointage} porte les deux horodatages exigés
par le cahier des charges --- l'heure d'arrivée et l'heure de départ --- ainsi
que l'attribut \texttt{source}, qui distingue le pointage automatique de la
saisie manuelle par l'administrateur.

\begin{coeur}
Trois questions tombent souvent sur ce diagramme.
\textbf{Pourquoi cinq paquetages ?} Pour la lisibilité : treize classes à plat
seraient illisibles, les paquetages regroupent par domaine fonctionnel.
\textbf{Pourquoi pas d'héritage ?} Parce qu'aucune entité n'est une
spécialisation d'une autre : les quatre types de demande ont des attributs et
des circuits différents, les regrouper sous une classe mère n'apporterait rien.
\textbf{Pourquoi un seul \texttt{@ManyToMany} ?} Parce qu'un seul lien est
réellement multiple des deux côtés : une demande d'agence concerne plusieurs
employés, et un employé peut figurer dans plusieurs demandes.
\end{coeur}

\begin{note}
\small Ne pas énumérer les treize classes : citer les cinq paquetages, désigner
\texttt{Utilisateur} au centre, et dire que les liens sont soit des clés
étrangères soit des relations JPA. Le jury cherche à savoir si vous comprenez
votre modèle, pas si vous l'avez appris par c\oe{}ur.
\end{note}

\subsection*{Slide 20 --- Modèle décisionnel en étoile}

J'ai comparé trois schémas --- étoile, flocon, constellation --- et retenu
l'étoile, plus simple et plus lisible pour un outil comme Power BI. La table de
faits \texttt{fait\_agence} regroupe sept indicateurs au grain de l'agence, avec
sa dimension \texttt{dim\_agence} ; une seconde table de faits,
\texttt{fait\_objectif}, descend au grain agence $\times$ gestionnaire, avec sa
dimension \texttt{dim\_gestionnaire}.

\begin{corrige}
La version d'origine disait \enquote{entourée de \textbf{quatre} dimensions :
agence, période, gestionnaire et type de client}. Le datamart en compte
\textbf{deux} : \texttt{dim\_agence} et \texttt{dim\_gestionnaire}. Les quatre
axes cités dans le rapport servent à \emph{justifier le choix de l'étoile}
(nombre réduit d'axes d'analyse) --- ce ne sont pas quatre tables de dimension.
Si le jury ouvre \texttt{etl/entrepot/}, il n'y trouvera que ces deux-là.
\end{corrige}

% =====================================================================
\section*{Partie 04 --- Volet décisionnel et réalisation (slides 21 à 29)}

\subsection*{Slide 22 --- Chaîne ETL en Python : cinq étapes}

La chaîne ETL a cinq étapes : la collecte des cinq tables sources, le nettoyage,
la transformation --- l'agrégation des indicateurs par agence ---, l'intégration
sur la clé \texttt{SK\_AGENCE}, et le contrôle de qualité. Cette chaîne
rejouable a produit \textbf{49 entités du réseau consolidées}.

\begin{corrige}
La version d'origine annonçait \enquote{45 agences consolidées avec leurs
7 indicateurs}. L'ETL produit \textbf{49 entités} : les 45 agences du réseau,
plus le siège, deux agences régionales et deux agences premium. C'est le chiffre
que porte le rapport et c'est celui qu'on lit dans la sortie du script. Le
nombre d'indicateurs dépend de la source : \textbf{sept} avec la table
\texttt{POINTAGE} (le taux de présence en fait partie), \textbf{six} sans elle.
\end{corrige}

\begin{note}
\small Si on vous demande \enquote{45 ou 49 ?} : \emph{45 agences
d'exploitation forment le réseau commercial ; l'extraction en remonte 49 parce
qu'elle inclut le siège et les structures régionales et premium.}
\end{note}

\subsection*{Slide 23 --- Datamart obtenu sur les données réelles}

Le datamart révèle une hétérogénéité entre agences --- BIZERTE et MGHIRA
dominent le portefeuille clients, CENTRALE la production de crédits --- que la
segmentation va exploiter.

\subsection*{Slide 24 --- Intégration du modèle dans Power BI}

Le modèle Power BI relie l'agence à ses employés, clients et objectifs, ce qui
permet un filtrage croisé : un clic recalcule instantanément tous les visuels du
rapport.

\subsection*{Slides 25 à 28 --- Tableaux de bord}

Quatre tableaux de bord Power BI sur données réelles : un résumé du réseau ---
30\,000 clients, 1,40 milliard de crédits, 87,9\,\% de présence ; l'analyse
clients --- portefeuille à 88\,\% particuliers, concentré sur le Grand Tunis ;
la performance commerciale --- production équilibrée, CENTRALE comme agence
phare ; et le suivi de présence --- bonne adoption du pointage automatique, avec
des agences à surveiller identifiées.

\begin{note}
\small Quatre diapositives : commenter un point par diapositive, sans revenir
en arrière.
\end{note}

\subsection*{Slide 29 --- Réalisation et environnement technique}

L'application repose sur Java 21, Jakarta EE, Hibernate et Oracle, déployée sur
WildFly. J'ai validé l'application avec dix scénarios de tests fonctionnels,
tous conformes.

\begin{note}
\small Les dix scénarios sont ceux du tableau du rapport (chapitre 4) :
authentification, rôles, agences, employés, clients, objectifs, pointage,
demandes, ETL, décisionnel. Les citer dans cet ordre en cas de question ---
ce sont ceux qui figurent dans le document remis au jury.
\end{note}

% =====================================================================
\section*{Partie 05 --- Intelligence artificielle et démonstration (slides 30 à 33)}

\subsection*{Slide 31 --- Zoom technique : segmentation par K-Means}

Le principe : regrouper les agences en profils homogènes à partir de leurs
indicateurs, sans étiquette préalable --- apprentissage non supervisé.

Le siège est écarté de l'analyse : avec 540 employés, c'est une structure
centrale et non une agence, et il écraserait toutes les distances. Les montants,
très asymétriques, passent au logarithme, puis toutes les variables sont
standardisées. J'ai retenu \textbf{k = 3} à l'aide de la méthode du coude et du
score de silhouette, puis comparé quatre algorithmes --- K-Means, hiérarchique,
GMM et DBSCAN.

Le modèle retenu est \textbf{K-Means}, avec un score de silhouette de
\textbf{0,604} et un indice de Calinski-Harabasz de \textbf{95,2}, les meilleurs
des quatre.

\begin{corrige}
La version d'origine annonçait \enquote{silhouette 0,736 et Davies-Bouldin
0,350}. Ces valeurs viennent de l'ancienne exécution sur un \textbf{jeu
synthétique} ; sur les données réelles, K-Means obtient \textbf{0,604} de
silhouette et \textbf{0,709} de Davies-Bouldin. Le tableau du rapport porte les
valeurs réelles : si vous citez 0,736, le jury verra la contradiction.
\end{corrige}

\begin{coeur}
\textbf{Attention à ne pas dire \enquote{les meilleurs sur les quatre
indices}.} C'est faux et vérifiable dans le tableau : K-Means a la meilleure
silhouette (0,604) et le meilleur Calinski-Harabasz (95,2), mais \textbf{DBSCAN
a un meilleur Davies-Bouldin} (0,393 contre 0,709). La bonne formulation :
\emph{\enquote{K-Means obtient la meilleure silhouette. DBSCAN a un
Davies-Bouldin plus favorable, mais il refuse de classer CENTRALE, qu'il traite
comme du bruit : je retiens le modèle qui classe toutes les agences, ce qui est
la condition pour s'en servir en pilotage.}}
\end{coeur}

\begin{note}
\small Sur \textbf{k = 3} : la silhouette est maximale à k = 2 (0,681), mais
k = 2 ne fait qu'isoler les entités sans activité. Le mieux est de le dire
soi-même : \emph{\enquote{k = 2 sépare trivialement les entités sans activité ;
j'ai retenu la meilleure solution à partir de k = 3, exploitable pour le
pilotage.}}
\end{note}

\subsection*{Slide 32 --- Résultats de la segmentation}

Trois profils émergent : \textbf{34 agences d'exploitation}, qui portent
l'activité du réseau ; \textbf{6 entités de production sans portefeuille
client} --- structures régionales et cas particuliers ; et \textbf{8 entités
sans activité enregistrée}. De quoi allouer les ressources et différencier les
objectifs selon le profil réel de chaque structure.

\begin{corrige}
La version d'origine annonçait \enquote{7 grandes agences, 16 moyennes,
22 petites}. C'est la répartition de l'ancien jeu synthétique. Sur les données
réelles, les trois groupes sont \textbf{34 / 6 / 8}, et ils ne se lisent pas en
\enquote{grandes, moyennes, petites} mais en \emph{nature d'activité} : celles
qui exploitent un portefeuille, celles qui produisent sans portefeuille, celles
sans activité enregistrée.
\end{corrige}

\subsection*{Slide 33 --- Démonstration}

Je vais maintenant vous faire une démonstration de l'application web, du tableau
de bord embarqué, et du rapport Power BI.

\begin{note}
\small Se taire pendant que la vidéo tourne, et commenter au fur et à mesure.
Avant la démonstration : vérifier qu'Oracle est démarré et que la base
pluggable est ouverte, sinon l'application renvoie \texttt{IJ000453}.
\end{note}

% =====================================================================
\section*{Partie 06 --- Conclusion et perspectives (slides 34 à 37)}

\subsection*{Slide 35 --- Conclusion}

Ce projet livre trois réalisations : une application web de gestion centralisée
du réseau, une chaîne décisionnelle complète avec quatre tableaux de bord
Power BI, et un volet intelligence artificielle de segmentation en trois profils
exploitables. Au final, un système complet, de la donnée opérationnelle à la
décision.

\subsection*{Slide 36 --- Perspectives}

Plusieurs perspectives : une sécurité renforcée --- le mot de passe est
aujourd'hui comparé en clair, il faudrait un condensat salé ---,
l'industrialisation de l'ETL, des analyses prédictives, une API REST, une
application mobile, et des alertes automatiques sous seuils critiques.

\subsection*{Slide 37 --- Clôture}

Je vous remercie pour votre attention, et je reste à votre disposition pour vos
questions.

% =====================================================================
\newpage
\section*{Récapitulatif des corrections}

\noindent\begin{tabularx}{\textwidth}{|>{\columncolor{btkpale}\bfseries}L{2.2cm}|L{5.3cm}|X|}
\hline
\textmd{\textbf{Slide}} & \textbf{Version d'origine} & \textbf{Version vérifiée}\\
\hline
6 & \enquote{issus du datamart produit par ma chaîne ETL} & issus des données réelles restituées par Power BI (l'ETL agrège un périmètre plus restreint)\\
\hline
20 & entourée de \textbf{quatre} dimensions & \textbf{deux} dimensions : \texttt{dim\_agence} et \texttt{dim\_gestionnaire}\\
\hline
22 & \textbf{45} agences consolidées & \textbf{49} entités du réseau consolidées\\
\hline
31 & silhouette \textbf{0,736}, Davies-Bouldin \textbf{0,350} & silhouette \textbf{0,604}, Davies-Bouldin \textbf{0,709} ; DBSCAN a un meilleur Davies-Bouldin\\
\hline
32 & \textbf{7} grandes / \textbf{16} moyennes / \textbf{22} petites & \textbf{34} d'exploitation / \textbf{6} de production sans portefeuille / \textbf{8} sans activité\\
\hline
\end{tabularx}

\vspace{0.8em}
\noindent\textbf{Chiffres vérifiés et inchangés :} 45 agences, 997 employés,
29\,942 clients, 105\,000 comptes, 1,40 milliard de crédits, 253 millions
d'épargne, 87,9\,\% de présence, 88\,\% de particuliers, création de la BTK le
25 février 1981, stage du 16 février au 16 juin 2026, dix scénarios de test.

\vspace{0.6em}
\begin{coeur}
Les trois chiffres à ne pas confondre le jour J : \textbf{45} agences dans le
réseau, \textbf{49} entités dans l'extraction, \textbf{48} retenues dans la
segmentation (le siège est écarté).
\end{coeur}

\end{document}
